\section{Dipol-Antenne}
\label{section:dipol}
\begin{frame}%STARTCONTENT
    \begin{figure}
        \DARCimage{0.85\linewidth}{589include}
        \caption{\scriptsize Darstellung einer Dipol-Antenne}
        \label{n_dipol}
    \end{figure}
    \begin{columns}
        \begin{column}{0.48\textwidth}
            \begin{itemize}
                \item In der Praxis wird häufig der Halbwellendipol verwendet
                \item Ist eine halbe Wellenlänge lang
            \end{itemize}
        \end{column}
        \begin{column}{0.48\textwidth}
            \pause
            \begin{exampleblock}{Beispiel:}
                \begin{itemize}
                    \item Wellenlänge von \qty{10}{\metre}
                    \item halbe Wellenlänge \qty{5}{\metre}
                    \item Jedes Teilstück des Dipols \qty{2,5}{\metre}
                \end{itemize}
            \end{exampleblock}
        \end{column}
    \end{columns}
\end{frame}

\begin{frame}
    \only<1>{
        \begin{PQuestion}{NG103}{Wie wird die dargestellte Antenne bezeichnet?}{Endgespeiste Antenne}
            {Groundplane-Antenne}
            {Dipol-Antenne}
            {Yagi-Uda-Antenne}
            {\DARCimage{1.0\linewidth}{589include}}\end{PQuestion}

    }
    \only<2>{
        \begin{PQuestion}{NG103}{Wie wird die dargestellte Antenne bezeichnet?}{Endgespeiste Antenne}
            {Groundplane-Antenne}
            {\textbf{\textcolor{DARCgreen}{Dipol-Antenne}}}
            {Yagi-Uda-Antenne}
            {\DARCimage{1.0\linewidth}{589include}}\end{PQuestion}

    }
\end{frame}

\begin{frame}
    \frametitle{Anpassung}
    Dipol-Antenne auf die gewünschte Frequenz bringen durch gleichmäßiges Kürzen oder Verlängern

    \begin{itemize}
        \item Zu hohe Resonanzfrequenz: Beide Seiten gleichmäßig verlängern
        \item Zu niedrige Resonanzfrequenz: Beide Seiten gleichmäßig verkürzen
    \end{itemize}
\end{frame}

\begin{frame}
    \only<1>{
        \begin{QQuestion}{NG304}{Ihre selbstgebaute Dipol-Antenne ist unterhalb der gewünschten Frequenz resonant. Welche Änderung können Sie vornehmen, um die Resonanz in den gewünschten Bereich zu bringen?}{Sendeleistung verringern}
            {Beide Enden gleichmäßig verlängern}
            {Sendeleistung erhöhen}
            {Beide Enden gleichmäßig kürzen}
        \end{QQuestion}

    }
    \only<2>{
        \begin{QQuestion}{NG304}{Ihre selbstgebaute Dipol-Antenne ist unterhalb der gewünschten Frequenz resonant. Welche Änderung können Sie vornehmen, um die Resonanz in den gewünschten Bereich zu bringen?}{Sendeleistung verringern}
            {Beide Enden gleichmäßig verlängern}
            {Sendeleistung erhöhen}
            {\textbf{\textcolor{DARCgreen}{Beide Enden gleichmäßig kürzen}}}
        \end{QQuestion}

    }
\end{frame}

\begin{frame}
    \only<1>{
        \begin{QQuestion}{NG305}{Ihre selbstgebaute Dipol-Antenne ist oberhalb der gewünschten Frequenz resonant. Welche Änderung können Sie vornehmen, um die Resonanz in den gewünschten Bereich zu bringen?}{Sendeleistung erhöhen}
            {Beide Enden gleichmäßig kürzen}
            {Beide Enden gleichmäßig verlängern}
            {Sendeleistung verringern}
        \end{QQuestion}

    }
    \only<2>{
        \begin{QQuestion}{NG305}{Ihre selbstgebaute Dipol-Antenne ist oberhalb der gewünschten Frequenz resonant. Welche Änderung können Sie vornehmen, um die Resonanz in den gewünschten Bereich zu bringen?}{Sendeleistung erhöhen}
            {Beide Enden gleichmäßig kürzen}
            {\textbf{\textcolor{DARCgreen}{Beide Enden gleichmäßig verlängern}}}
            {Sendeleistung verringern}
        \end{QQuestion}

    }
\end{frame}%ENDCONTENT
